\documentclass[10pt,a4paper]{amsart}
\usepackage[margin=19mm]{geometry}
\usepackage{amsmath,amssymb,mathtools}
\usepackage[scr=rsfs,cal=euler]{mathalfa}
\usepackage{xfrac}
\usepackage[unicode,hidelinks,bookmarks=false]{hyperref}
\newcommand{\ie}{i.e.\ }
\newcommand{\cf}{cf.\ }
\newcommand{\resp}{respectively\ }
\newcommand{\Subsection}{\subsection}
\providecommand{\nobreakdash}{\nobreak\hbox{-}}
\setlength{\emergencystretch}{2em}
\multlinegap0pt
\usepackage{mathtools}
\usepackage{xcolor}
\usepackage{amssymb}
\usepackage{dsfont}
\usepackage{tikz}
\newtheorem{lemma}{Lemma}
\newtheorem{prop}{Proposition}
\usepackage{cleveref}
\crefname{lemma}{Lemma}{Lemmas}
\crefname{prop}{Proposition}{Propositions}
\DeclareMathOperator{\Aut}{Aut}
\title{A simple proof of the fundamental\\theorem of Galois theory}
\author{Martin Brandenburg}
\date{Source: arXiv:2211.07508v2 (CC BY 4.0)}
\begin{document}
\maketitle

\noindent\emph{Source and licence.} A simple proof of the fundamental\\theorem of Galois theory. Version arXiv:2211.07508v2 (CC BY 4.0), distributed by arXiv under the Creative Commons Attribution 4.0 International licence, http://creativecommons.org/licenses/by/4.0/, as recorded in the arXiv licence metadata for this exact version and shown on the abstract page https://arxiv.org/abs/2211.07508v2. Full licence notice: \texttt{licenses/arxiv-2211.07508-CC-BY-4.0.txt}.\par\smallskip

\noindent\textit{Editorial scope.} Counted unit: the article's prop subgroupclass (source0.tex lines 227--230). The article's complete original proof of it is delivered with it (source0.tex lines 232--242, one \texttt{proof} environment of 11 lines). No mathematics is written, repaired or re-derived: the statement and the proof are the article's own lines, in the article's own notation, and no retained line is substituted or re-flowed.  Retained uncounted as the source-local context the proof needs: lines 199--201.  Nothing was omitted from any retained span, so the package carries no omission record.  No retained line contains a citation, so the wrapper carries no bibliography and no bibliography entry.  No retained line contains a figure environment, an image-inclusion command or drawing code, so no image is delivered and the package declares no asset dependency.  Editorial formatting only, in addition: an AMS \texttt{amsart} preamble that restates the article's own declarations and macros used by the retained lines, namely the environments \texttt{lemma}, \texttt{prop} on separate counters, together with the standard packages those lines need.  The retained environments print in this document as Lemma 1, Proposition 1 in order of appearance, whereas the article prints the same items with the numbering of its own full document.  The complete native source of the article as distributed in the arXiv e-print for this exact version is bundled unchanged as \texttt{sources/133459/source0.tex}.
\begin{lemma} \label{fieldunion}
A field cannot be written as the union of finitely many proper subfields.
\end{lemma}

\begin{prop} \label{subgroupclass}
Let $L/K$ be an extension and $H \subseteq \Aut_K(L)$ be a finite subgroup. Then
\[H = \Aut_{L^H}(L).\]
\end{prop}

\begin{proof}
Let $\tau \in \Aut_{L^H}(L)$. We need to prove $\tau \in H$. First, we claim that
\[L = \bigcup_{\sigma \in H} \{\sigma = \tau\},\]
where $\{\sigma = \tau\}$ is a short notation for the subfield $\{a \in L : \sigma(a) = \tau(a)\}$, the equalizer of $\sigma,\tau$. Let $a \in L$. In order to use the assumption that $\tau$ fixes elements of $L^H$, we need to somehow come up with elements of $L^H$. Consider the polynomial
\[p \coloneqq \prod_{\sigma \in H} \bigl(T - \sigma(a)\bigr) \in L[T].\]
Of course, this is only well-defined since $H$ is finite, and we have $p(a)=0$. The natural action of $H$ on $L$ extends to an action on $L[T]$. The polynomial $p$ is clearly fixed by this action since the action just permutes the linear factors. Hence, we have
\[p \in L[T]^H = L^H[T].\]
Thus, $\tilde{\tau} : L[T] \to L[T]$ fixes $p$, i.e.\ $\tilde{\tau}(p)=p$. Since $\tau(a)$ is a root of $\tilde{\tau}(p)=p$, there is some $\sigma \in H$ with $\sigma(a)=\tau(a)$, so that $a \in \{\sigma = \tau\}$. This proves our claim.

Now, each equalizer $\{\sigma = \tau\}$ is a subfield of $L$. Thus, \cref{fieldunion} implies that there is some $\sigma \in H$ with $L = \{\sigma = \tau\}$, which just means $\sigma = \tau$.
\end{proof}

\end{document}
